\chapter{绪论}\label{chap:introduction}

\section{研究背景}
随着在线教学、远程会议与跨地域协作的常态化，团队在讨论过程中的“即时表达”需求不断增强。相较于传统文档协作，白板能够更自然地承载草图、流程图与推导过程，尤其适合用于快速构思与讲解。在实际使用中，用户往往希望在尽量低的门槛下迅速进入协作状态，即“打开即可用、创建即可分享、加入即可同步”。

现有白板产品在功能完整性与使用成本之间通常存在取舍：平台化产品往往引入账号体系、团队空间与复杂权限配置；而极简产品虽然上手快，但在多人同步、网络波动与跨端一致性方面可能出现体验问题。基于上述背景，设计并实现一个聚焦核心协作链路、以房间化方式降低使用成本的在线白板系统具有现实意义。

\section{研究意义}
本课题的意义主要体现在以下方面：
\begin{itemize}
  \item \textbf{工程实践价值}：综合应用现代前端框架、图形渲染与实时通信技术，完成从需求到交付的系统实现与验证闭环。
  \item \textbf{协作体验价值}：以房间号与邀请链接为核心入口，降低进入协作的操作成本，提升“创建—分享—加入”的效率。
  \item \textbf{可扩展性价值}：通过清晰的模块边界与同步策略，为后续扩展更多图形组件、增强一致性与持久化能力提供基础。
\end{itemize}

\section{国内外研究与现状}
在线白板的发展经历了从单机绘图工具到 Web 化协作平台的演进。产品形态上，一类系统强调极简与快速进入，适合轻量协作；另一类系统更强调团队化与流程化能力，适用于复杂项目管理。技术路线方面，前端多采用基于 Canvas 或 WebGL 的图形渲染方案，协作层多基于 WebSocket 类实时通信实现状态广播与同步；在冲突处理上，常见策略包括按时间戳的最后写入优先、服务端权威裁决以及更复杂的 CRDT/OT 等一致性算法。

在教学与工程讨论场景中，快速创建房间并共享给同伴、在网络波动下仍能保持基本可用，是用户体验的关键。QuickBoard 选择以“轻量房间协作 + 必要绘图工具 + 数学公式识别闭环”为目标，不追求平台化功能，但重点保证核心链路的稳定与快速。

\section{本文主要工作}
围绕 QuickBoard 在线白板系统，本文完成的主要工作包括：
\begin{itemize}
  \item 设计并实现基于房间的白板协作入口，支持创建房间、加入房间与最近房间管理；
  \item 实现基于 Node.js 与 Socket.IO 的实时协作通信，并采用服务端权威裁决与增量同步策略提升一致性体验；
  \item 集成 OCR 服务实现数学公式识别，并在前端实现渲染与二次编辑流程；
  \item 通过单元测试与构建验证等方式对系统进行验证，并在交互与渲染侧进行针对性优化以提升稳定性与响应速度。
\end{itemize}

\section{论文结构}
本文结构安排如下：第 \ref{chap:tech} 章介绍系统实现所涉及的关键技术；第 \ref{chap:req} 章进行需求分析；第 \ref{chap:design} 章给出系统总体架构与模块设计；第 \ref{chap:impl} 章阐述关键实现；第 \ref{chap:eval} 章对系统进行测试与评估；第 \ref{chap:conclusion} 章给出总结与展望。

